\documentclass[11pt]{article}
\usepackage[margin=1in]{geometry}
\usepackage{enumitem}
\usepackage{hyperref}
\usepackage[T1]{fontenc}
\setlist[itemize]{leftmargin=1.4em, itemsep=0.25em, topsep=0.3em}

\title{EPA Air Emissions Programs (POLL\_RPT\_COMBINED\_EMISSIONS)}
\author{}
\date{}

\begin{document}
\maketitle

\noindent The \texttt{POLL\_RPT\_COMBINED\_EMISSIONS} dataset combines facility-level emissions from four
EPA programs. This overview states only facts confirmed against EPA documentation or verified directly
in the data (sources at the end).

\section*{The four programs}
\begin{itemize}
  \item \textbf{NEI --- National Emissions Inventory.} EPA's National Emissions Inventory is ``released
  every three years'' (triennial) and is a comprehensive estimate of air emissions of ``criteria
  pollutants, criteria precursors, and hazardous air pollutants'' from air-emission sources.
  \item \textbf{CAMD --- Clean Air Markets / Acid Rain Program.} The Acid Rain Program requires major
  reductions of sulfur dioxide and nitrogen oxide emissions from electric utilities through a ``cap and
  trade'' approach, with affected sources required to install continuous emissions monitors that report
  measured (not estimated) emissions.
  \item \textbf{GHGRP --- Greenhouse Gas Reporting Program.} EPA reports that approximately 8,000
  facilities are required to report their greenhouse-gas emissions annually.
  \item \textbf{TRI --- Toxics Release Inventory.} Established by the Emergency Planning and Community
  Right-to-Know Act (EPCRA, 1986; Section 313 of EPCRA established the TRI). EPA reports that the TRI
  program currently covers 810 individually listed chemicals and 34 chemical categories. TRI is a
  separate statute from the Clean Air Act.
\end{itemize}

\section*{Fields}
Columns of \texttt{POLL\_RPT\_COMBINED\_EMISSIONS.csv} (verified against the file header):
\begin{itemize}
  \item \texttt{REPORTING\_YEAR} --- calendar year of the emission report.
  \item \texttt{REGISTRY\_ID} --- Facility Registry Service (FRS) ID; the join key to ICIS-Air.
  \item \texttt{PGM\_SYS\_ACRNM} --- source program: \texttt{EIS} (NEI), \texttt{E-GGRT} (GHGRP),
  \texttt{TRIS} (TRI), \texttt{CAMDBS} (CAMD).
  \item \texttt{PGM\_SYS\_ID} --- program-specific facility identifier.
  \item \texttt{POLLUTANT\_NAME} --- name of the reported pollutant.
  \item \texttt{ANNUAL\_EMISSION} --- emission value for the facility, pollutant, and year.
  \item \texttt{UNIT\_OF\_MEASURE} --- \texttt{Pounds} or \texttt{MTCO2e} (metric tons
  CO\textsubscript{2}-equivalent).
  \item \texttt{NEI\_TYPE} --- pollutant category: \texttt{CAP} (criteria air pollutant), \texttt{GHG}
  (greenhouse gas), \texttt{HAP} (hazardous air pollutant), \texttt{OTH}, \texttt{PFAS}.
  \item \texttt{NEI\_HAP\_VOC\_FLAG} --- flag identifying a hazardous-air-pollutant volatile organic
  compound.
\end{itemize}

\section*{Considerations}
\begin{itemize}
  \item \textbf{Mixed units.} Greenhouse-gas values are in metric tons CO\textsubscript{2}-equivalent;
  other values are in pounds. The two units cannot be summed without conversion.
  \item \textbf{Different reporting universes and frequencies.} NEI is triennial while CAMD, GHGRP, and
  TRI report annually, and each program has its own reporting criteria, so a facility may appear in one
  program but not another and coverage differs by program and year.
\end{itemize}

\section*{Sources}
\begin{itemize}
  \item EPA, \emph{National Emissions Inventory (NEI)},
  \url{https://www.epa.gov/air-emissions-inventories/national-emissions-inventory-nei}
  \item EPA, \emph{Greenhouse Gas Reporting Program}, \url{https://www.epa.gov/ghgreporting}
  \item EPA, \emph{What is the Toxics Release Inventory?},
  \url{https://www.epa.gov/toxics-release-inventory-tri-program/what-toxics-release-inventory}
  \item EPA, \emph{AFS Data Download} documentation (April 2014) --- Acid Rain Program description.
  \item Direct verification against \texttt{data/raw/POLL\_RPT\_COMBINED\_EMISSIONS.csv}.
\end{itemize}

\end{document}
